\documentclass[10pt,a4paper]{amsart}
\usepackage[margin=19mm]{geometry}
\usepackage{amsmath,amssymb,mathtools}
\usepackage[scr=rsfs,cal=euler]{mathalfa}
\usepackage{xfrac}
\usepackage[unicode,hidelinks,bookmarks=false]{hyperref}
\newcommand{\ie}{i.e.\ }
\newcommand{\cf}{cf.\ }
\newcommand{\resp}{respectively\ }
\newcommand{\Subsection}{\subsection}
\providecommand{\nobreakdash}{\nobreak\hbox{-}}
\setlength{\emergencystretch}{2em}
\multlinegap0pt
\usepackage{caption}
\usepackage{booktabs}
\usepackage{amssymb}
\usepackage{tikz}
\usepackage[normalem]{ulem}
\usepackage{xcolor}
\usepackage{mathtools}
\usepackage[shortlabels]{enumitem}
\newtheorem{lem}{Lemma}
\newcommand{\Ga}{\Gamma}
\newcommand{\bcs}{\prescript{}{\pd}{\#}}
\newcommand{\cd}{\cdots}
\newcommand{\e}{\epsilon}
\newcommand{\ga}{\gamma}	
\newcommand{\pd}{\partial}
\title{On a conjecture of Almgren II: area-minimizing submanifolds with fractal singular sets on almost any manifold}
\author{Zhenhua Liu}
\date{Source: arXiv:2512.24859v1 (CC BY 4.0)}
\begin{document}
\maketitle

\noindent\emph{Source and licence.} On a conjecture of Almgren II: area-minimizing submanifolds with fractal singular sets on almost any manifold. Version arXiv:2512.24859v1 (CC BY 4.0), distributed by arXiv under the Creative Commons Attribution 4.0 International licence, http://creativecommons.org/licenses/by/4.0/, as recorded in the arXiv licence metadata for this exact version and shown on the abstract page https://arxiv.org/abs/2512.24859v1. Full licence notice: \texttt{licenses/arxiv-2512.24859-CC-BY-4.0.txt}.\par\smallskip

\noindent\textit{Editorial scope.} Counted unit: the article's lem lembcsd (source0.tex lines 1338--1340). The article's complete original proof of it is delivered with it (source0.tex lines 1341--1343, one \texttt{proof} environment of 3 lines). No mathematics is written, repaired or re-derived: the statement and the proof are the article's own lines, in the article's own notation, and no retained line is substituted or re-flowed.  Retained uncounted as the source-local context the proof needs: lines 765--765.  Nothing was omitted from any retained span, so the package carries no omission record.  No retained line contains a citation, so the wrapper carries no bibliography and no bibliography entry.  No retained line contains a figure environment, an image-inclusion command or drawing code, so no image is delivered and the package declares no asset dependency.  Editorial formatting only, in addition: an AMS \texttt{amsart} preamble that restates the article's own declarations and macros used by the retained lines, namely the environments \texttt{lem} on separate counters, together with the standard packages those lines need.  The retained environments print in this document as Lemma 1 in order of appearance, whereas the article prints the same items with the numbering of its own full document.  The complete native source of the article as distributed in the arXiv e-print for this exact version is bundled unchanged as \texttt{sources/191936/source0.tex}.
	\subsection{Construction of coordinate chart $U$}\label{seccht}

	\begin{lem}\label{lembcsd}
		If $W,W'$ are two smooth domains of $M,$ and $\ga$ is a smooth curve from $w\in \pd W$ to $w'\in \pd W'$ such that $\ga$ meets $W,W'$ only at $w,w'$, is transverse to $\pd W,\pd W',$ and is orientation reversing. Then for any neighborhood $U(\ga)$ of $\ga,$ there exists a smooth embedding $\Ga$ of $W\bcs W'$ into $M$ and $\Ga(W\bcs W')\setminus (W\cup W')$ is diffeomorphic to a neck $B_1^{d+c-1}\times[-1,1]$ contained in $U(\ga).$
	\end{lem}

	\begin{proof}
		The proof is similar to the construction in Section \ref{seccht}. First take a coordinate chart $U$ with coordinate labels $(x_1,\cd,x_{d+c})$, so that $W$ is the half space $x_{d+c}\le -\frac{1}{2}$, $W'$ is the half space $x_{d+c}\ge \frac{1}{2}$ and $\ga$ is the line segment from $(0,\cd,0,-\frac{1}{2})$ and $(0,\cd,0,\frac{1}{2}).$ Glue a neck $B_\e^{d-1}(\{0\}^{d-1})\times\{0\}^c\times [-\frac{1}{2},\frac{1}{2}]$ for $\e$ small and smooth the corners. We are done.
	\end{proof}

\end{document}
